% sec-00: front matter.  Abstract, the nothing-sent notice, the reading path.
\section*{Abstract}
\addcontentsline{toc}{section}{Abstract}

Labor Day is a federal holiday and a full closure of the New York Stock Exchange
and Nasdaq \cite{opmholidays2026,nyseholidays2026}. No counterparty this program
deals with can act, which makes it the one day of the block on which everything
can be prepared without interruption.

\textbf{Nothing in this packet or in the directory that holds it has been sent,
entered, filed, or agreed, and nothing may be before the next open session.} Four
letters are written and held. Six orders are sized and priced and exist in no
broker system. A nine-folder data room is assembled. A twenty-two question
diligence bank is built, and four of its questions are marked as having no
answer.

The single approval this day asks for is therefore different in kind from the
other four in the block. It is not "send this." It is "authorize the release
list, so that the next session opens with the queue already approved."

\section{Why a Closed Day Is in the Schedule at All}

An email that lands in a program officer's inbox on a federal holiday arrives at
the bottom of a Tuesday stack, behind everything that accumulated while the
office was closed. An order entered against a closed session either rejects or
rests until the open and fills at a price formed in the first minutes of a
session nobody has seen, which is worse because it looks like a normal fill.

Both failures are avoidable by writing today and releasing on the next open
session. That is the whole design.

There is a second reason, and it is the better one. A letter written under time
pressure is a worse letter. The most consequential item in this five-day block is
a Pre-Request for Designation asking which FDA center leads the review of a
system that is a drug, a robotic platform, and an advisory software component at
once \cite{fdacombination2026}. Every regulatory assumption downstream depends on
its answer. Writing it between two other things on a busy session would produce
the version that fits the window rather than the version that is right.

\subsection{How to read this}

A regulator reads \S1 and \S5 and nothing else: the system and its stop
conditions. Those two are the only sections attached to the classification
letter.

An investor's analyst under a confidentiality agreement reads \S3 and \S5: the
data room index and the diligence bank, including the four questions with no
answer.

The chief executive reads all seven, because it is his approval the day asks for,
and \S2 carries it.
